\documentclass[11pt]{article}
\usepackage[a4paper,margin=1in]{geometry}
\usepackage{amsmath,amssymb,amsthm,mathtools}
\usepackage{graphicx}
\usepackage{microtype}
\usepackage{xcolor}
\usepackage{enumitem}
\usepackage[colorlinks=true,linkcolor=blue!55!black,citecolor=blue!55!black,urlcolor=blue!55!black]{hyperref}

\newtheorem{theorem}{Theorem}
\newtheorem{proposition}{Proposition}
\newtheorem{lemma}{Lemma}
\newtheorem{remark}{Remark}
\newcommand{\R}{\mathbb R}
\newcommand{\I}{\mathcal I}
\newcommand{\norm}[1]{\left\lVert#1\right\rVert}

\title{A sharp Fourier gap for Z\"ust's coefficient-kernel question}
\author{Automated mathematics run \texttt{fa\_banach\_001}\\Lane 14, candidate partial result}
\date{11 August 2026}

\begin{document}
\maketitle

\begin{center}
\fbox{\parbox{0.91\textwidth}{\textbf{Status.} Candidate substantial partial result, likely valid pending expert review.  We compute the exact Hessian at the north pole, prove that every open-hemisphere profile is stationary, and prove strict nonlinear local maximality on the Fourier-normal slice.  The source's global inequality and equality characterization remain open.}}
\end{center}

\section{Source question and reformulation}

The source is Roger Z\"ust, \emph{The Riemannian hemisphere is almost calibrated in the injective hull of its boundary}, arXiv:2104.04498v2 (submitted 2021, revised 2025) \cite{Zust}.  Its final Question~2, on PDF page~48, asks whether
\[
 \norm{p(f)}_1:=\int_0^\pi\int_\alpha^\pi p_{\alpha,\beta}(f)\,d\beta\,d\alpha
 \le \frac{\pi^2}{2}
\]
for every $f\in E(\mathbf S^1)\setminus\mathbf S^1$, with equality exactly when $f$ belongs to the embedded open hemisphere $\mathbf S^2_+\setminus\mathbf S^1$.

\begin{figure}[ht]
  \centering
  \includegraphics[width=0.98\textwidth]{figures/open_problem_crop.png}
  \caption{The definition and Question~2 in the source PDF, page~48.}
\end{figure}

In the source's lifted model, $f:\R\to(0,\pi)$ is $1$-Lipschitz and
$f(\alpha+\pi)+f(\alpha)=\pi$.  Put
\[
 a(\alpha)=\cos f(\alpha).
\]
Then $a(\alpha+\pi)=-a(\alpha)$ and
\begin{equation}\label{eq:slope}
 |a'(\alpha)|^2\le 1-a(\alpha)^2\quad\text{a.e.}
\end{equation}
Conversely, these conditions with $|a|<1$ recover an interior point by
$f=\arccos a$.  Writing $c=\cos(\beta-\alpha)$ and
$s=\sin(\beta-\alpha)$, the coefficient is
\begin{equation}\label{eq:p}
 p_{\alpha,\beta}(a)=
 \frac{s^2-a(\alpha)^2-a(\beta)^2+2c\,a(\alpha)a(\beta)}
 {s^2(1-a(\alpha)^2)(1-a(\beta)^2)}.
\end{equation}
The embedded hemisphere consists precisely of first harmonics
$a(\alpha)=A\cos\alpha+B\sin\alpha$ with $A^2+B^2<1$.

\section{Proof intuition}

At the north pole $a=0$, the kernel is identically one, so the functional is
$\pi^2/2$.  The numerator of \eqref{eq:p} is quadratic in $a$, but the denominator contains all higher nonlinearities.  The quadratic part is nevertheless exactly diagonal in the odd Fourier basis.  Mode~$1$ is neutral because it moves inside the hemisphere; every mode $n\ge3$ loses exactly $(\pi^2/2)(n-1)$ times its squared amplitude.

The nonlinear terms retain this sign near zero.  An exact identity separates the quadratic kernel from a favorable negative quartic term.  The only indefinite remainder is the quadratic kernel multiplied by a factor of order $\norm a_\infty^2$.  A fractional-difference-quotient estimate controls its absolute integral by the same Fourier energy as the Hessian.  On the slice with no first harmonic, the spectral gap absorbs the remainder.

This does not globalize automatically: the denominators couple Fourier modes, and removing the first harmonic is only a linear normal gauge at the north pole.  A separate weighted Hilbert-transform calculation does show that every hemisphere profile is stationary in all directions, strengthening the local geometric picture but not supplying global coercivity.

\section{Exact Hessian and local theorem}

For an antiperiodic real function, use the normalization
\begin{equation}\label{eq:fourier}
 a(\alpha)=\sum_{\substack{n\ge1\\n\ \mathrm{odd}}}
 \bigl(A_n\cos(n\alpha)+B_n\sin(n\alpha)\bigr).
\end{equation}
Define
\begin{equation}\label{eq:q}
 q_{\alpha,\beta}(a)=
 \frac{2c\,a(\alpha)a(\beta)-c^2(a(\alpha)^2+a(\beta)^2)}{s^2},
 \qquad
 Q[a]=\frac12\int_0^\pi\int_0^\pi q_{\alpha,\beta}(a)\,d\beta\,d\alpha.
\end{equation}

\begin{theorem}[Sharp north-pole Hessian]\label{thm:hessian}
For every smooth antiperiodic $a$,
\begin{equation}\label{eq:spectrum}
 Q[a]= -\frac{\pi^2}{2}
 \sum_{\substack{n\ge3\\n\ \mathrm{odd}}}(n-1)(A_n^2+B_n^2).
\end{equation}
Consequently, the Hessian of $\I(a):=\frac12\int_0^\pi\int_0^\pi p_{\alpha,\beta}(a)$ at $a=0$ is negative semidefinite, and its kernel is exactly $\operatorname{span}\{\cos\alpha,\sin\alpha\}$, the tangent plane to the embedded hemisphere.
\end{theorem}

\begin{proof}
The bilinear polarization of $Q$ is translation invariant on the antiperiodic circle, so distinct odd Fourier modes are orthogonal.  For either sine or cosine mode $n$ with squared amplitude $A_n^2+B_n^2$, direct integration in the base variable gives
\[
 Q_n=\frac\pi2(A_n^2+B_n^2)
 \int_0^\pi\frac{\cos t\cos(nt)-\cos^2t}{\sin^2t}\,dt.
\]
Write $n=2k+1$.  The finite identity
\[
 \frac{\cos(nt)-\cos t}{\sin t}=-2\sum_{j=1}^k\sin(2jt)
\]
and the elementary integral
$\int_0^\pi\cot t\,\sin(2jt)\,dt=\pi$ give
\[
 \int_0^\pi\frac{\cos t\cos(nt)-\cos^2t}{\sin^2t}\,dt
 =-\pi(n-1).
\]
Summing the mutually orthogonal modes proves \eqref{eq:spectrum}.  Expanding \eqref{eq:p} at zero shows that $Q$ is its quadratic term; there is no linear term.  The kernel statement follows immediately.
\end{proof}

\begin{lemma}[Absolute quadratic-kernel bound]\label{lem:absolute}
There is an absolute constant $C$ such that every Lipschitz antiperiodic $a$ satisfies
\begin{equation}\label{eq:absolute}
 \int_0^\pi\int_0^\pi |q_{\alpha,\beta}(a)|\,d\beta\,d\alpha
 \le C\sum_{\substack{n\ge1\\n\ \mathrm{odd}}}n(A_n^2+B_n^2).
\end{equation}
\end{lemma}

\begin{proof}
For $t=|\beta-\alpha|\le\pi/2$, the numerator in \eqref{eq:q} can be written
\[
 q=-\frac{c^2(a(\alpha)-a(\beta))^2}{s^2}
   +\frac{2c}{1+c}a(\alpha)a(\beta).
\]
For $t\ge\pi/2$, instead use
\[
 q=-\frac{c^2(a(\alpha)+a(\beta))^2}{s^2}
   +\frac{2c}{1-c}a(\alpha)a(\beta).
\]
The second terms are bounded by a constant times $a(\alpha)^2+a(\beta)^2$.  In the first long-distance term, antiperiodicity gives
$a(\alpha)+a(\beta)=a(\beta)-a(\alpha+\pi)$, and its circular separation is $\pi-t\le\pi/2$.  Thus the two singular terms are bounded by the usual fractional difference quotient on the $2\pi$-circle.  The standard Fourier identity for that quotient is a constant multiple of
$\sum |n|(A_n^2+B_n^2)$; Parseval controls the remaining $L^2$ terms.  This proves \eqref{eq:absolute}.
\end{proof}

\begin{theorem}[Nonlinear strict local maximum on the Fourier-normal slice]\label{thm:local}
There is $\varepsilon_0>0$ such that the following holds.  Let
$f\in E(\mathbf S^1)\setminus\mathbf S^1$, put $a=\cos f$, and suppose
\[
 \norm a_\infty\le\varepsilon_0,
 \qquad A_1=B_1=0.
\]
Then
\begin{equation}\label{eq:local}
 \norm{p(f)}_1\le\frac{\pi^2}{2}
 -\frac{\pi^2}{4}
 \sum_{\substack{n\ge3\\n\ \mathrm{odd}}}(n-1)(A_n^2+B_n^2).
\end{equation}
Equality holds only for the north pole $f\equiv\pi/2$.
\end{theorem}

\begin{proof}
Put $a=a(\alpha)$, $b=a(\beta)$ and
$D=(1-a^2)(1-b^2)$.  Direct common-denominator algebra gives the exact identity
\begin{equation}\label{eq:identity}
 p_{\alpha,\beta}(a)-1=\frac{q_{\alpha,\beta}(a)-a^2b^2}{D}.
\end{equation}
Therefore
\begin{align}
 \I(a)-\I(0)
 &=Q[a]+\frac12\iint q_{\alpha,\beta}(a)(D^{-1}-1)
   -\frac12\iint a^2b^2D^{-1}.\label{eq:remainder}
\end{align}
The last term is nonpositive.  If $\norm a_\infty\le1/2$, then
$0\le D^{-1}-1\le C_0\norm a_\infty^2$.  By Lemma~\ref{lem:absolute}, the absolute value of the middle term in \eqref{eq:remainder} is at most
\[
 C_1\norm a_\infty^2
 \sum_{\substack{n\ge1\\n\ \mathrm{odd}}}n(A_n^2+B_n^2).
\]
Since $A_1=B_1=0$ and $n\le\frac32(n-1)$ for odd $n\ge3$, this is at most
$C_2\norm a_\infty^2 E(a)$, where
$E(a)=\sum_{n\ge3,\,n\text{ odd}}(n-1)(A_n^2+B_n^2)$.
Theorem~\ref{thm:hessian} says $Q[a]=-(\pi^2/2)E(a)$.  Choose
$\varepsilon_0$ so that $C_2\varepsilon_0^2\le\pi^2/4$.  Then
\eqref{eq:remainder} gives \eqref{eq:local}.  Equality forces $E(a)=0$; the missing first mode then gives $a=0$, hence $f\equiv\pi/2$.
\end{proof}

\section{A hemisphere-wide stationarity identity}

\begin{proposition}[Every hemisphere profile is stationary]\label{prop:stationary}
Let $h(\alpha)=r\cos(\alpha-\tau)$ with $|r|<1$.  For every smooth antiperiodic variation $v$,
\[
 \left.\frac d{dt}\right|_{t=0}\I(h+tv)=0.
\]
\end{proposition}

\begin{proof}
Rotation lets us take $\tau=0$.  Set $k=(1-r^2)^{1/2}$ and
$A_x=1-r^2\cos^2x$.  The identity
\[
 \cos^2x+\cos^2y-2\cos(y-x)\cos x\cos y=\sin^2(y-x)
\]
gives $p_{x,y}(h)=k^2/(A_xA_y)$ and hence $\I(h)=\pi^2/2$.  Differentiating the symmetrized integral reduces the coefficient of $v(x)$ to
\[
 \frac{2r}{A_x}\left[-\operatorname{pv}\!\int_0^\pi
 \frac{\sin y}{\sin(y-x)A_y}\,dy
 \frac{k^2\cos x}{A_x}\int_0^\pi\frac{dy}{A_y}\right].
\]
The elementary weighted Hilbert-transform identities
\[
 \int_0^\pi\frac{dy}{A_y}=\frac\pi k,
 \qquad
 \operatorname{pv}\!\int_0^\pi
 \frac{\sin y}{\sin(y-x)A_y}\,dy=\frac{\pi k\cos x}{A_x}
\]
make this coefficient zero.  For completeness, the second identity follows by writing
$\sin y/\sin(y-x)=\cos x+\sin x\cot(y-x)$ and inserting the absolutely convergent Fourier expansion
\[
 \frac1{1-r^2\cos^2y}=\frac1k\left(1+2\sum_{m\ge1}
 \left(\frac{1-k}{1+k}\right)^m\cos(2my)\right).
\]
The paired first variation is locally integrable across $x=y$; the principal-value notation only records the symmetrization used above.
\end{proof}

\section{Verification, limitations, and upgrade audit}

The script \texttt{code/verify\_kernel.py} checks the exact nonlinear identity on 1,000 random nonsingular pairs and compares midpoint quadrature for modes $1,3,5,7,9,11$ with \eqref{eq:spectrum}.  Independent constrained discretizations with 12, 20, and 32 nodes and dozens of starts found no value above $\pi^2/2$.  These computations are sanity checks, not proof.

Eight focused routes were pursued: (1) the cosine-coordinate/Gram reduction; (2) finite Lipschitz optimization; (3) a de Sitter/Lorentz reformulation; (4) the exact Fourier Hessian; (5) nonlinear remainder absorption; (6) the full-hemisphere stationarity identity; (7) a pointwise product domination, disproved by the family $a=0.2\cos(3\alpha)$; and (8) a global Fourier/Morse--Bott normal form.  The last route is obstructed because the nonlinear denominators couple all modes and the first-harmonic projection is not a global admissible gauge.  Neither stationarity nor the north-pole gap currently gives a uniformly coercive transverse Hessian around the full hemisphere.

The cheap run indexes and bounded exact-phrase, title, author, citation, and keyword searches found no later resolution of Question~2.  The present novelty confidence is moderate pending specialist review.  The global inequality, behavior far from the hemisphere, and global equality characterization remain open.

\section*{Human-review recommendation}

Treat Theorems~\ref{thm:hessian} and \ref{thm:local} and Proposition~\ref{prop:stationary} as a likely valid substantial partial.  The highest-value checks are the normalization in the full-square symmetrization, the fractional $H^{1/2}$ estimate at both singular diagonals, and differentiation under the integral in Proposition~\ref{prop:stationary}.  A promising next mathematical step would be a weighted transverse Hessian around $r\cos(\alpha-\tau)$ that is uniform for $r$ in compact subsets of $(-1,1)$.

\begin{thebibliography}{9}
\bibitem{Zust}
R.~Z\"ust,
\emph{The Riemannian hemisphere is almost calibrated in the injective hull of its boundary},
arXiv:2104.04498v2, 2021; revised 2025.
\end{thebibliography}

\end{document}
